\section{Ecuaciones diferenciales ordinarias (EDO) lineales y soluciones exactas}

\subsection{Solución analítica de EDO lineales simples}

Antes de discutir las EDO del modelo de difusión, el conferenciante parte de la EDO lineal más simple para construir intuición. Considere una EDO escalar:
$$\frac{dx}{dt} = p(t) \cdot x(t)$$

donde $p(t)$ es una función conocida del tiempo. Suponga que $x(t) > 0$, divida ambos lados por $x(t)$:
$$\frac{1}{x(t)} \frac{dx}{dt} = p(t)$$

El lado izquierdo es exactamente $\frac{d}{dt} \ln x(t)$, por lo tanto:
$$\ln x(t) - \ln x(s) = \int_s^t p(\tau)\, d\tau$$

\begin{importantbox}{Solución exacta de EDO lineales}
$$x(t) = x(s) \cdot \exp\left(\int_s^t p(\tau)\, d\tau\right)$$
Esta solución es\textbf{analítica}: no requiere ninguna aproximación por discretización. Para EDO lineales, siempre podemos obtener una solución exacta mediante integración. Este hecho es la piedra angular clave para las derivaciones posteriores.
\end{importantbox}

\subsection{EDO semilineales y el método del factor integrante}

La PF-ODE de los modelos de difusión no es puramente lineal, sino\textbf{semilineal}, con la forma:
$$\frac{d\mathbf{x}}{dt} = p(t) \cdot \mathbf{x}(t) + q(t)$$

donde $p(t) \cdot \mathbf{x}$ es la parte lineal y $q(t)$ es el término de desplazamiento no lineal (en los modelos de difusión, $q(t)$ contiene la salida de la red predictora del ruido).

Para resolver una EDO semilineal, se introduce un\textbf{factor integrante}:
$$\mu(t) = \exp\left(-\int_0^t p(\tau)\, d\tau\right)$$

$\mu(t)$ es la "solución inversa" de la ecuación homogénea $\frac{dx}{dt} = p(t) \cdot x$. Multiplique ambos lados de la EDO por $\mu(t)$:
$$\mu(t) \frac{dx}{dt} - \mu(t) p(t) x(t) = \mu(t) q(t)$$

Utilizando la regla del producto para derivación, el lado izquierdo se simplifica a $\frac{d}{dt}\left[\mu(t) x(t)\right]$; al integrar, se obtiene:

\begin{importantbox}{Solución exacta de EDO semilineales (Variación de constantes)}
$$x(t) = \frac{\mu(s)}{\mu(t)} x(s) + \frac{1}{\mu(t)} \int_s^t \mu(\tau) q(\tau)\, d\tau$$
\begin{itemize}
    \item Primer término: Solución exacta de la parte lineal, depende únicamente del valor inicial $x(s)$ y puede calcularse analíticamente.
    \item Segundo término: Integral de la parte no lineal, involucra a $q(\tau)$ (que contiene la salida de una red neuronal) y generalmente requiere aproximación numérica.
\end{itemize}
Esto es la forma básica del\textbf{integrador exponencial}: se trata la parte lineal con precisión exacta y solo se realiza una aproximación numérica para la parte no lineal.
\end{importantbox}

\begin{knowledgebox}{¿Por qué se llama "integrador exponencial"?}
Debido a que la solución exacta de la parte lineal incluye la función exponencial $\exp(\cdot)$. El método de Euler tradicional realizaría una aproximación sobre toda la EDO (lineal + no lineal), mientras que el integrador exponencial solo aproxima la parte no lineal. En casos donde la parte lineal domina (como en la PF-ODE de los modelos de difusión), esta división mejora significativamente la precisión.
\end{knowledgebox}

\subsection{Solución exacta con término de desplazamiento constante}

Como ejercicio, considere el caso con una constante $q$ (independiente de $t$):
$$\frac{dx}{dt} = p(t) \cdot x(t) + c$$

Al sustituir la fórmula de la solución general y simplificar, se puede obtener una expresión más compacta. El conferenciante llevó a los estudiantes en clase a derivar este proceso; la idea central fue sacar $c$ fuera del signo integral y luego simplificar utilizando las propiedades de $\mu(t)$.

\subsection{Resumen de la sección}

Las EDO lineales tienen soluciones analíticas; las EDO semilineales pueden resolverse exactamente en su parte lineal mediante el método del factor integrante, dejando solo la parte no lineal que requiere aproximación numérica. Este es el fundamento matemático de DPM-Solver.
